\begin{figure}[htbp]
\centering
\begin{tikzpicture}[font=\small,>=Stealth]
\node[align=center,font=\small\bfseries] at (0,2) {Old presentation};
\node[align=center,font=\small\bfseries] at (6,2) {Add $a\simeq_{\mathcal P}e;e$};
\node at (0,1.5) {$|f|=|e;e|$};
\node at (6,1.5) {$|a|=|f|=|e;e|$};
\draw[rounded corners,draw=pathblue,thick] (-1.65,-0.2) rectangle (-0.2,0.8);
\draw[rounded corners,draw=pathblue,thick] (0.2,-0.2) rectangle (1.65,0.8);
\node at (-0.93,0.5) {$[f]$}; \node at (-0.93,0.1) {parity $1$};
\node at (0.93,0.5) {$[e;e]$}; \node at (0.93,0.1) {parity $0$};
\draw[rounded corners,draw=pathorange,thick] (4.05,-0.2) rectangle (7.95,0.8);
\node at (4.9,0.5) {$[f]$}; \node at (6.8,0.5) {$[a]=[e;e]$};
\node at (6,0.05) {one-letter $a,f$: same code};
\draw[->,thick] (1.95,0.55) -- (3.75,0.55)
  node[midway,above] {inclusion};
\draw[<-,pathorange,thick] (1.95,-0.55) -- (3.75,-0.55)
  node[midway,below] {$a\mapsto e;e$};
\node[align=center] at (0,-1.25) {Distinct classes;\\continuous parity separates them};
\node[align=center] at (6,-1.25) {Distinct classes;\\identical open neighborhoods};
\end{tikzpicture}
\caption{A primitive-to-word substitution preserves the presented algebra
but can change the observable topology (Proposition~\ref{prop:redundant-generator-observable}).
Adding the defining rule gives inverse scoped substitutions. Inclusion is
continuous on the observable quotient; elimination is discontinuous because
it sends inseparable classes to classes separated by parity. Boxes indicate
separation properties, not asserted open two-point subspaces.}
\FigureAltText{Old classes f and e;e are separated by length parity. After
adding a, its class equals e;e, while the one-letter representatives a and f share an observable code. Eliminating a
sends inseparable classes to separated old classes.}
\label{fig:redundant-generator}
\end{figure}
